\question[5]
Su x86 (little-endian): (a) l'intero \lstinline|0x41424344| come appare come sequenza di 4 byte in memoria (dall'indirizzo più basso)? (b) La sequenza di byte \lstinline|10 00 00 00| quale intero rappresenta? (c) Perché \lstinline|p32()|/\lstinline|p64()| sono comode?
\begin{solution}
(a) \lstinline|44 43 42 41| (byte meno significativo per primo).\\
(b) \lstinline|0x00000010| $=$ \textbf{16}.\\
(c) Impacchettano automaticamente un intero nei byte in little-endian (p32 su 4 byte, p64 su 8), evitando errori manuali nell'ordine dei byte.
\end{solution}
